\chapter[Temples, Ancestors, Invisible Order]{Temples, Ancestors, and Invisible Order}
\section{A Cracked Turtle Shell}
Pick something up---not in your hand, but behind museum glass, where it rests on dark velvet.

A turtle's lower shell, roughly palm-sized, has the colour of old bone and a web of fine cracks. Look closely: most cracks resemble \mbox{卜}, bu, "divination", a vertical stroke and a shorter sideways stroke, like little signposts. Beside them run tiny engraved characters, thin and hard, their turns retaining the knife point's sharpness.

The label identifies a Shang oracle bone over three millennia old, excavated at Yinxu in Anyang, Henan. Over a hundred thousand inscribed pieces have been found. Before them, knowledge of Shang relied largely on later books. Their emergence was like Shang people handing up a hundred thousand notes about daily concerns.

Those concerns may surprise you: rain, toothache caused by ancestors, tomorrow's campaign, the queen's unborn child's sex, tonight's sacrifice. Every shell is a solemn question addressed not to living people but ancestors, Di, and invisible powers of mountains, rivers, wind, rain.

Do not immediately say superstition. The easy word often stops thinking for us. Try something harder: imagine a world where unseen powers are not optional beliefs but airlike assumptions. Then examine what temples, incense, divination, ancestral tablets actually do in lives. Returning afterward to today, you may recognise unexpected acquaintances.

\section{An Evening in the Great Settlement Shang}
Enter a scene.

Time: dusk around thirty-two centuries ago. Place: the late Shang capital its inhabitants called the Great Settlement Shang, around today's Anyang.

Humid heat rises from the Huan River. A palace courtyard is meticulously swept and matted. Well-dressed people work beside stacks of turtle shells and cattle shoulder blades, not casually collected bones but distant tribute requiring soaking, scraping, polishing.

Tonight's charge asks whether the royal army should depart tomorrow.

A divination official, whose practice the Shang called zhen and whom we may call a diviner, raises a shell in both hands, formally addresses an unseen audience, and states the question. Then a red-hot hardwood rod touches a prepared hollow on its reverse.

Crack.

A slight sharp sound, and a \mbox{卜}-shaped crack opens on the front. Everyone leans closer. That sound is the answer. Direction, length, angle convey ancestral and divine advice requiring expert reading, with final judgement belonging to the king. Inscriptions show Shang kings often personally declared good or bad omens.

Afterward the diviner engraves the date, questioner, charge, crack, royal judgement, and eventual verification. Filed bones resemble archived meeting minutes.

Watching from the corner, your first feeling may be seriousness rather than magic. No joking or perfunctory gestures; every step has rules governing questioner, bone, heating, reading, engraving. Less a fraudulent spectacle than a rigorous administrative meeting, with over half the participants invisible.

\section{Were They Merely Superstitious?}
Set out the conflict without rushing to answers.

A familiar account says ancient scientific ignorance invented gods for thunder and weather, while helplessness sought divine reassurance. Science would eventually retire these practices. Through this lens, Shang elites solemnly surrounding cracked shell collectively misunderstood, paying what today's slang calls a stupidity tax.

Several facts resist that account.

First, Shang ranked among East Asia's most technically advanced societies. Hundreds-of-kilograms bronzes with intricate patterns still impress artisans; cities had rammed-earth walls and palace foundations; writing was systematic. These people were not stupid. Calling divination ignorance merely supposes bronze-casting minds could not distinguish cracks from weather forecasts, replacing explanation with ancient foolishness.

Second, why would the most powerful, best-resourced king put reassurance for the naive at the state's centre? Why establish elaborate turtle procurement, shell preparation, charges, engraving, archives, more regulated than many modern expense claims? Mere foolishness rarely supports such heavy institutions.

Third, consider yourself. A lucky exam pen you refuse to replace? A fixed pre-match action? Reposted good-luck pictures despite knowing their impotence? Most people have such habits. Scientific understanding and ritual needs coexist peacefully in one mind. Are we judging ancient people by another standard?

The real question does not ask whether they believed in gods, inaccessible minds we cannot inspect. It asks \textbf{what visible practices---burial, sacrifice, divination, ancestor offerings---actually accomplished in people's lives}.

Why would a society spend so much maintaining invisible order? Harder and more interesting than labelling superstition, this takes the rest of our chapter step by step.

\section{The King Hears Ancestors; Who Hears Him?}
Position changes what one divination reveals. Hear several voices.

\textbf{First, the king.}

Divination grounds authority, not a hobby. Not anyone can address ancestors or Di, the Shang's supreme ruler. Dead ancestors remain active, blessing and harming the living. Whose ancestors rank highest and connect best to Di? The king's. Ultimate authority over war, crops, capital moves belongs naturally to him, not through volume but his family's exclusive highest connection.

Give this reasoning fully without endorsing it. King and subjects imagine a hierarchical house: Di highest, royal ancestors between, living people below. Acting on great matters without consultation resembles spending all company funds without board approval, reckless rather than brave. Divination is decision's proper procedure, not evasion, checking mortal judgement against larger order. A king bypassing it would seem unreasonable.

\textbf{Second, ordinary people.}

Inscriptions mostly concern royalty because engraving and storing bones were court luxuries. Ordinary voices leave no notes, but village shrines, sacrificial pits near homes, and grave pottery offer clues. Their unseen authorities were closer: grandparents, local land deities, hearth and doorway guardians. Concerns centred on feverish children, difficult cattle births, yearly rain, not campaigns.

This pattern continued in China: village earth-god temples, city-god shrines, merchants' wealth gods, boat people's Mazu. Gods subcontract jurisdictions, protecting particular sectors and places. Natural powers of weather and water, local supervisors, occupational and family protectors form a dense, tangible administrative network. A travelling child often prompts offerings to a road deity rather than supreme heaven, because a particular journey falls under a particular office.

An important asymmetry: royal ancestors affect the world, ordinary ancestors their own household. Invisible hierarchy aligns precisely with visible rank. The level at which one may burn incense itself declares status.

\textbf{Third, the offerings themselves.}

This heaviest voice must not vanish. Near Yinxu royal tombs and palaces, sacrificial pits hold groups of human bones, many Qiang war captives. Inscriptions also record captive sacrifice to ancestors. Shang cosmology may find this logical: grand ancestors deserve grand offerings. Logic cannot speak for victims. These bones belonged to named people with homes and waiting loved ones, whose names no oracle bone preserved.

Remember this voice because explaining ritual's social operation easily becomes cold admiration of smooth machinery without seeing those caught within it. Explaining an institution and defending it are different. This sentence will recur throughout the book; here it receives its first solemn statement.

Before proceeding, mention evidence far older than Shang: \textbf{burial}. Doing something for dead companions may predate writing, farming, temples. At Zhoukoudian's Upper Cave, people tens of thousands of years ago buried the dead deep inside with red hematite powder, blood's colour. Their thoughts remain unknowable; motives are conjecture. Yet they could have discarded bodies like rubbish, and did not. They spent effort. Three thousand years later in Shang, this effort grew immense. Fu Hao, King Wu Ding's consort who led armies and rites, received nearly two thousand bronze, jade, bone, and other grave objects, buried royally near the palace to continue participating in family order. A handful of Upper Cave red powder and her treasures mark one thread: refusal to let death end relationships. Burial may be our earliest invisible order.

\textbf{Pull back to other civilisations' versions of the same drama.}

Around thirty-eight centuries ago, Babylonian Hammurabi's famous code began not with provisions but divine commission: gods entrusted justice to their chosen champion of the people. Authority's source precedes law's content. Kingship needs a celestial letter of introduction.

Across the Mediterranean, Greek city-states consulted Delphi before war, colonisation, legislation. Apollo's priestess supplied ambiguous, solemn answers. Oracles issued no commands, yet no city dared act without asking.

Across the Pacific, sixteenth-century Aztecs believed, their tradition's claim rather than a historical conclusion, that daily solar struggle required hearts and blood or the world would darken. Priests accordingly organised large sacrifices. Spanish conquerors arriving centuries later recorded barbarism and committed another great violence in their own faith's name. Each civilisation believed it held the universe's manual.

Together, a pattern appears: most ancient societies connect who decides with who reaches the powers above. Not coincidence, but structure to examine.

\section{Thinking Bridge: Dismantle a Ritual}
Unfamiliar ancestral rites, prayers, ceremonies commonly trigger only belief or superstition. Both settings are too coarse for analysis. Our portable tool, \textbf{ritual analysis}, asks four questions.

\textbf{First: Who attends, and where do they stand?}

Ritual rarely allows casual arrival. Presiding, kneeling first, watching outside, exclusion: positions map social relationships. Royal crack-reading also means others cannot read with equivalent authority. At Qingming graveside visits, who carries photographs, burns paper, or may bow without incense reveals age and closeness. At weddings, flag ceremonies, school openings, first watch positions rather than content.

\textbf{Second: Why this time?}

Ritual rarely occurs whenever convenient. It fixes seasonal joints: before sowing, after harvest, year's beginning and end, full moon, seventh or forty-ninth day after death, anniversaries. Time flows continuously; rituals plant stakes marking different days. A community's calendar ranks what matters.

\textbf{Third: How must actions be done, and what if they go wrong?}

Correct bows, paper quantities, words, sequences matter. Consequences of mistakes may be unclear, yet everyone senses wrongness, invalidity, unease. This is a core component: \textbf{correct procedure itself reassures}. Modern law strikingly resembles it: a missing signature can invalidate truthful contractual content. Ritual is pre-legal signature culture.

\textbf{Fourth: What does it synchronise?}

Most important. One time, place, movement, direction tunes many minds together. Harvest sacrifice turns scattered households into our village; state funerals make strangers our nation; graduation makes classmates our cohort. Sociologists say ritual produces the group itself. Blood and maps do not automatically create groups; repeated shared acts refresh them, ritual the most efficient refresh key.

These questions leave belief's dead end for observable, comparable ground. Applied equally to ancient and modern, they measure Yinxu divination and your school's flag raising with one ruler. We will test it at the end.

\section{Confucius: Sacrifice as If Present}
Read primary material. Shang notes lay underground for millennia, but several centuries after Shang a Spring and Autumn teacher made a subtle observation.

The Analects' "Eight Rows" records:

\begin{quote}
He sacrificed as if they were present, and to spirits as if spirits were present. The Master said: "If I do not participate in sacrifice, it is as though I have not sacrificed."
\end{quote}
Meaning: treat ancestors and spirits as present; without personal participation, Confucius considers sacrifice unperformed.

Weigh ru, "as if". He asserts neither presence nor absence, but \textbf{the attitude of their presence during the rite}. In the same chapter, another statement makes his attitude clearer, cited as "Yong Ye": respect spirits while keeping a distance, and one may be called wise.

Together these passages occupy a narrow middle position: neither eager affirmation nor destructive dismissal; meticulous rites with existence gently bracketed. The focus lies with living sincerity, disposition, attendance, not spirits. Substitution fails because personal presence is what ritual purchases; proxy shopping is invalid.

A later reminder in "Learning" comes from disciple Zengzi: attentive funerals and remembrance of distant ancestors deepen popular virtue. Notice causation: ancestral rites concretely benefit \textbf{living hearts}. Whether ancestors receive offerings goes undiscussed; whether living people become kinder receives confident affirmation.

Twenty-five centuries ago, insiders already suspected ritual's delivery address might not be heaven. Never imagine ancient traditions as one solid, unquestioningly devout voice. Debate, doubt, and minds redirecting attention have always lived within them.

\section{One Incense Stick, Three Accounts}
Place Shang courtyard divination, Confucius's as-if presence, and Qingming paper ashes beneath one question: \textbf{why do people do these things?}

Separate what happens, timed formal acts toward unseen recipients; why, the following dispute; participants' accounts, ancestors need offerings or sincerity of presence; and our eventual evaluations.

\textbf{Explanation 1: They truly believe.}

The most direct and least dismissible answer. Kings believe ancestors speak through cracks; farmers believe dragon kings control rain; Aztec priests believe the sun needs feeding. Honest attention takes testimony seriously, without casting people as fools or actors. Extensive texts and objects indicate sincerity. Yet belief cannot explain this ritual rather than another, power's proximity to gods, or Roman statesman and actual augur Cicero's simultaneous service and systematic scepticism about prediction. Belief explains sincerity, not its specific shape.

\textbf{Explanation 2: Ritual is society's cement.}

Ignore internal belief; inspect work done. Gather people, beat time, arrange rank, recharge belonging. Spring rites coordinate fieldwork; funerals renew family relations; coronations publicly confirm rulers. This explains compulsory attendance by unbelievers, governmental investment, reluctance to cancel even apparently empty ceremonies. But translating everything into function misses real awe and grief. Describing someone weeping over ashes as executing social integration may be partly true yet chilling. Function explains shape, not tears' warmth.

\textbf{Explanation 3: Ritual manages uncertainty.}

Anthropologists studying fishing communities found few rites among calm coastal fishers, many strict ones among crews facing dangerous offshore seas. Same sea, different risks. Shang toothache and war invite divination because uncontrollable; lunch menus do not appear in inscriptions because too controllable to trouble ancestors. This predicts ritual concentration in risk and uncertainty and explains scientifically sceptical good-luck reposting. Yet it cannot explain ritual's \textbf{publicness}. Private reassurance could happen under a blanket. Why entire villages, elaborate gatherings, generations keeping one date? Individual comfort cannot explain collective persistence.

Each grasps truth and misses something. Not compromise: understanding their different jobs gives three lamps for future rituals instead of blindness from one.

\section[Further Challenge: Society Worships Itself]{A Deeper Challenge: Society Worships Itself}
Our strongest theory comes from Émile Durkheim's 1912 The Elementary Forms of Religious Life. Studying Aboriginal Australian totemic institutions, he reached a startling conclusion: \textbf{a group worshipping its gods is actually worshipping itself}.

How? Diverse religions divide things into \textbf{sacred} and \textbf{profane}, everyday. Sacred objects cannot be casually touched, places casually entered, days casually worked. Why forbid touching a wooden totem? It represents the group, not itself. Awe before wood is awe before the community, too large and abstract to feel without a visible vessel.

Several puzzles open. Desecrating flags, graves, sacred objects provokes fury exceeding material value because it affronts us. Coordinated ritual moves people to tears by dissolving individuality into something larger. Durkheim accepts the felt immense power as real but locates its source in assembled people, not heaven. Gatherings genuinely charge one another.

This rearranges earlier accounts. Royal divination switches on communal order; Qingming affirms continued family; the fed Aztec sun illuminates a people's place in the world.

Yet test the boundary. What of solitary midnight prayer, hermits apart from society, a mother's unwitnessed bedside vow to any god? No crowd, no collective charging, sometimes opposition to communal expectations. Durkheim illuminates public squares better than bedside whispers. \textbf{Religious belief, social function, and personal experience are three real, entangled things; none can swallow another completely.} Belief denying function is arrogant; function dissolving belief cold; experience denying both naive. Carry this into nearly every later discussion of faith.

\section{Lucky Fish, Exams, and Flag Raising}
This ancient question lives daily at school.

Before entrance exams, Confucian temples reach peak attendance: devout parents and students seeking merely a lucky omen. Online koi reposts promise success. Most know scores will not change, but just-in-case reassurance moves fingers. Athletes similarly repeat songs, avoid court lines, preserve team entrance sequences for years. They know these gestures do not determine results. Yet control itself has value: perform fully controllable acts before an uncontrollable outcome. Exactly uncertainty theory's prediction: ritual density follows risk.

At group level, memorial days lower flags, sound sirens, and synchronise silence; Mondays align thousands toward one school flag. Apply four questions. Positions: teachers, pupils, central flag bearers. Time: the week's threshold. Mistakes: unspecified consequences but talking feels wrong. Synchrony: thousands of separate concerns tune together. Structurally related to Yinxu's courtyard, radically changed in content, the skeleton survives astonishingly.

Ancestors moved rather than departed. Qingming offers homeward grave visits or virtual candles and messages. Japanese households present food at Buddhist altars and welcome ancestral returns during Obon. Mexican Day of the Dead displays favourite foods and photographs, believing, within the tradition, that departed relatives reunite. Elsewhere people simply remember inwardly. Which is more civilised may be the wrong question: all translate the ancient need to maintain relationships with those gone.

A real dispute deserves hearing. Each Qingming, paper burning provokes charges of pollution, forest fires, purchased reassurance, answered by tradition and emotional necessity. Give defenders their full case without declaring victory. They buy not paper but a remaining channel for caring. Loved ones depart; care's momentum needs somewhere to go. Fire sends undeliverable concern, seemingly providing warmth once more. Calling that ignorance dismisses another's most tender place. Yet opponents also have concrete reasons: burned forests, emergency-room injuries, cleaners' predawn ashes impose costs on others. When strangers pay for my reassurance, it ceases to be private. Full hearings reshape the question from right-or-wrong burning to ways of holding grief without transferring its costs. Asking that makes our tools worthwhile.

\section{Your Turn: If Every Ritual Vanished}
Return to the cracked shell.

An ancient king entrusted military departure to heated turtle bone. We neither can nor need return, but its world persists in changed clothes outside exams, inside phones, on playgrounds.

Imagine every ritual vanishing: no mourning halls, weddings, grave visits, flags, silence, birthdays, exam prayers, or funerary procedure, bodies simply disposed of efficiently. Would the world become more rational or lose a supporting corner?

Answer with reasons after recounting tools. Four questions reveal group-making through position, time, procedure, synchrony. Three explanations show belief, cement, reassurance, indispensable yet conflicting. Durkheim's gatherings genuinely produce belonging despite his overconfident source claim. Cicero and Confucius show participation need not require full belief, nor unbelief escape ritual.

Is ritual a crutch to discard when strong enough, or a skeleton without which we lose human shape? Does a civilisation without sacrifice, remembrance, celebration finally awaken, or forget itself?

No standard answer. Next Monday's anthem may briefly recall that ancient dusk: people around shell, breath held for a small crack. Millennia, science, everything separates you, yet the act may be closer than expected. How close is yours to answer.
